\section{Three unequal frequencies and all their regularized jets}
\label{xid:sec:triple-dilation}

The current incoming report treats three distinct translates and asks
for three unequal integer frequencies \cite{IncomingIdentities}.
Fourier orthogonality leaves a lattice of rank two in this problem.
Finite congruence filters nevertheless give an explicit finite family
of colored Tornheim kernels.  Local subtraction then converts the
spectral identity into an all-index Stieltjes identity in the ambient
coordinate.  The classical use of colored Tornheim sums and their
polylogarithmic reductions is credited to the literature
\cite{ZhaoColored}; the formulas here address the unequal-frequency
and fixed-coordinate question from the supplied research.

Let $p_1,p_2,p_3$ be positive integers and $a_1,a_2,a_3$ be real
shifts, considered modulo one.  Write
\[
 S_i=\{x\in\mathbb R/\mathbb Z:p_i x+a_i\in\mathbb Z\},
 \qquad e(t)=e^{2\pi i t}.
\]
Throughout this section assume that $S_1,S_2,S_3$ are pairwise
disjoint.  Equivalently, for every $i\ne j$,
\begin{equation}\label{xid:eq:triple-separation}
 p_j a_i-p_i a_j\notin\gcd(p_i,p_j)\mathbb Z.
\end{equation}
At each grid point the cutoff is the right ambient coordinate
$t=x-x_0$; the cutoffs at different grid points can tend to zero
independently.

\subsection{Weighted colored kernels and finite character filters}

For positive integers $a,b$ and nontrivial unit complex numbers $X,Y$
define initially in an absolutely convergent half-space
\begin{equation}\label{xid:eq:weighted-T}
 T_{a,b}(A,B,C;X,Y)
 =\sum_{m,n\ge1}\frac{X^mY^n}{m^A n^B(am+bn)^C}.
\end{equation}
The positive real denominators use real logarithms.  For $\Re C>0$
the entire continuation in $A,B$ is given by
\begin{equation}\label{xid:eq:weighted-T-mellin}
 T_{a,b}(A,B,C;X,Y)
 =\frac1{\Gamma(C)}\int_0^\infty t^{C-1}
        \Li_A(Xe^{-at})\Li_B(Ye^{-bt})\,dt.
\end{equation}
In fact this is jointly entire in all three orders after continuation.
To prove the statement, subtract any finite Taylor polynomial at
$t=0$ from the product of polylogarithms.  It is holomorphic there,
locally uniformly in $A,B$, because $X,Y\ne1$.  The Taylor integrals
have poles only at $C=0,-1,\ldots$, each canceled by $1/\Gamma(C)$.
The integral over $[1,\infty)$ converges locally normally for all
orders because its integrand decays exponentially.  Increasing the
Taylor order proves entire continuation and licenses every fixed
mixed order derivative.

These weighted kernels introduce no extra class of functions beyond
ordinary colored Tornheim kernels.  Choose any roots $\xi^a=X$ and
$\eta^b=Y$.  Applying two finite root filters gives
\begin{equation}\label{xid:eq:weighted-T-reduction}
 T_{a,b}(A,B,C;X,Y)
 =a^{A-1}b^{B-1}
 \sum_{r=0}^{a-1}\sum_{s=0}^{b-1}
 T_{1,1}(A,B,C;\xi e(r/a),\eta e(s/b)).
\end{equation}
Different choices of roots merely permute the summands.  None of
these twists is one, since $X,Y\ne1$.  The proof first changes
variables $u=am,v=bn$ in the absolutely convergent series, and then
uses $a^{-1}\sum_r e(ru/a)=\mathbf1_{a\mid u}$ and its $b$ analogue.
Analytic continuation proves the identity at all orders.

For each $k\in\{1,2,3\}$ let $i<j$ be the remaining indices, and
for $h=0,\ldots,p_k-1$ put
\begin{align}
 X_{ik,h}&=e\left(a_i-\frac{p_i a_k}{p_k}
                              +\frac{hp_i}{p_k}\right),\notag\\
 X_{jk,h}&=e\left(a_j-\frac{p_j a_k}{p_k}
                              +\frac{hp_j}{p_k}\right),\label{xid:eq:triple-twists}\\
 \vartheta_k&=\frac\pi2(\lambda_i+\lambda_j-\lambda_k).
 \notag
\end{align}
Every twist in \eqref{xid:eq:triple-twists} is nontrivial.  Indeed,
$X_{ik,h}=1$ for some integer $h$ is equivalent to
$p_k a_i-p_i a_k\in\gcd(p_i,p_k)\mathbb Z$, which contradicts
\eqref{xid:eq:triple-separation}.

\begin{theorem}[Finite spectral formula for three unequal frequencies]
\label{xid:thm:triple-spectral}
For $\Re\lambda_i>0$,
\begin{equation}\label{xid:eq:triple-K}
 K(\boldsymbol\lambda;\boldsymbol p,\boldsymbol a)
 =\int_0^1\prod_{i=1}^3
       \zeta(1-\lambda_i,\{p_i x+a_i\})\,dx
 =\frac{\prod_i\Gamma(\lambda_i)}{(2\pi)^{\sum_i\lambda_i}}
       \mathcal S(\boldsymbol\lambda),
\end{equation}
where the finite sum is
\begin{align}
 \mathcal S(\boldsymbol\lambda)
 =\sum_{k=1}^3p_k^{\lambda_k-1}\sum_{h=0}^{p_k-1}
 \bigl\{&e^{-i\vartheta_k}
 T_{p_i,p_j}(\lambda_i,\lambda_j,\lambda_k;
                         X_{ik,h},X_{jk,h})\notag\\
 &+e^{i\vartheta_k}
 T_{p_i,p_j}(\lambda_i,\lambda_j,\lambda_k;
                         X_{ik,h}^{-1},X_{jk,h}^{-1})\bigr\}.
 \label{xid:eq:triple-cones}
\end{align}
The right side is a meromorphic continuation to all three order
variables.  It has $2(p_1+p_2+p_3)$ weighted-kernel terms before
symmetry or character simplifications.
\end{theorem}

\begin{proof}
Start with $\Re\lambda_i>1$, where each Hurwitz Fourier series is
absolutely convergent.  For a nonzero frequency $n$, its coefficient
is
\[
 \frac{\Gamma(\lambda_i)}{(2\pi|n|)^{\lambda_i}}
      e^{-i\pi\lambda_i\operatorname{sgn}(n)/2}.
\]
After dilation and translation, integration retains precisely
$p_1n_1+p_2n_2+p_3n_3=0$, with all $n_i\ne0$.
There are six disjoint sign cones.  In the cone with $n_i=m>0$,
$n_j=n>0$ and $n_k<0$, one has
\[
 -n_k=\frac{p_i m+p_j n}{p_k}\in\mathbb N.
\]
The denominator contributes $p_k^{\lambda_k}$ and the phase of
the Fourier coefficient is $e^{-i\vartheta_k}$.  Insert
\[
 \mathbf1_{p_k\mid p_i m+p_j n}
   =\frac1{p_k}\sum_{h=0}^{p_k-1}
                  e\left(\frac{h(p_i m+p_j n)}{p_k}\right).
\]
Combining this filter with the translation phases produces exactly
the first term in braces in \eqref{xid:eq:triple-cones}, including the
factor $p_k^{\lambda_k-1}$.  Reversing all frequency signs gives
the second term, after reindexing the finite filter by $h\mapsto-h$.
Each lattice point occurs once.

The original integral is jointly holomorphic for $\Re\lambda_i>0$:
near a grid point only one factor is singular, with a locally
integrable majorant $t^{\sigma-1}$ on any compact parameter set.
Order derivatives add only logarithmic powers.  Entire continuation
of the colored kernels makes the right side holomorphic in this
half-space as well.  The identity theorem proves the asserted
extension, and the Gamma factors give the meromorphic continuation.
\end{proof}

There is also a longer formula directly in the known equal-frequency
kernel.  The ordinary Hurwitz multiplication law gives
\begin{equation}\label{xid:eq:covering-triple}
 K(\boldsymbol\lambda;\boldsymbol p,\boldsymbol a)
 =\prod_i p_i^{\lambda_i-1}
 \sum_{h_1=0}^{p_1-1}\sum_{h_2=0}^{p_2-1}\sum_{h_3=0}^{p_3-1}
 K\left(\boldsymbol\lambda;(1,1,1),
        \left(\frac{a_i+h_i}{p_i}\right)_{i=1}^3\right).
\end{equation}
All three translated endpoints in every summand are distinct.  This
provides an alternative finite reduction and confirms that the rank-two
lattice does not require a new, unbounded family of kernels.

\subsection{All Stieltjes and argument jets in the ambient coordinate}

Set $\mathscr G(\lambda,y)=\zeta(1-\lambda,y)+1/\lambda$ as before,
and define the pair kernels by
\[
 C_{ij}(\lambda_i,\lambda_j)
 =\int_0^1\zeta(1-\lambda_i,\{p_ix+a_i\})
                \zeta(1-\lambda_j,\{p_jx+a_j\})\,dx,
\]
with their ordinary separated-grid analytic continuation.
They have the finite Gamma--polylogarithm expression
\eqref{xid:eq:pair-spectral}. All shift derivatives below are taken within
a local chamber of separated real shifts. Define
\begin{equation}\label{xid:eq:triple-regular-spectral}
 \mathcal L(\boldsymbol\lambda)
 =K(\boldsymbol\lambda)+\sum_{\substack{\{i,j,k\}=\{1,2,3\}\\i<j}}
       \frac{C_{ij}(\lambda_i,\lambda_j)}{\lambda_k}
       +\frac1{\lambda_1\lambda_2\lambda_3}.
\end{equation}
In an ordinary convergent domain this is the integral of the product
of the three $\mathscr G$ functions.  The terms with just one Hurwitz
factor vanish because its integral over a period is zero.

For $\boldsymbol r=(r_1,r_2,r_3)\in\mathbb N_0^3$ let
$\partial_{\boldsymbol a}^{\boldsymbol r}=\prod_i\partial_{a_i}^{r_i}$.
For $x_0\in S_i$ and $j\ne i$ put
$\eta_{ij}(x_0)=\{p_jx_0+a_j\}\in(0,1)$.  Define the finite local
numerator
\begin{align}
 B_i^{\boldsymbol r}(\boldsymbol\lambda)
 ={}&(-1)^{r_i}(1-\lambda_i)_{r_i}
          p_i^{\lambda_i-r_i-1}\notag\\
 &\quad\cdot\sum_{x_0\in S_i}[t^{r_i}]
   \prod_{j\ne i}\partial_y^{r_j}
          \mathscr G(\lambda_j,\eta_{ij}(x_0)+p_jt).
 \label{xid:eq:triple-local-B}
\end{align}
It is holomorphic near the spectral origin.  It uses a finite Taylor
polynomial and values of Stieltjes functions and their argument
derivatives at nonsingular points.

\begin{theorem}[All-index triple finite-part identity]
\label{xid:thm:triple-FP}
The function
\begin{equation}\label{xid:eq:triple-holomorphic-J}
 \mathcal J_{\boldsymbol r}(\boldsymbol\lambda)
 =\partial_{\boldsymbol a}^{\boldsymbol r}
       \mathcal L(\boldsymbol\lambda)
       -\sum_{i=1}^3\frac{B_i^{\boldsymbol r}(\boldsymbol\lambda)}{\lambda_i}
\end{equation}
is holomorphic in a sufficiently small neighborhood of the origin.
For all $m_1,m_2,m_3\ge0$,
\begin{align}
 &\FP_x\int_0^1\prod_{i=1}^3
       \gamma_{m_i}^{(r_i)}(\{p_ix+a_i\})\,dx\notag\\
 &\qquad=\left(\prod_i m_i!\right)
       [\lambda_1^{m_1}\lambda_2^{m_2}\lambda_3^{m_3}]
          \mathcal J_{\boldsymbol r}(\boldsymbol\lambda).
 \label{xid:eq:triple-FP-value}
\end{align}
Together with \eqref{xid:eq:triple-cones}, this is an explicit finite
colored-Tornheim, Gamma, and local Stieltjes coefficient formula for
every choice of the six indices.
\end{theorem}

\begin{proof}
Begin with $\Re\lambda_i>r_i+1$. Argument differentiation and
integration then commute, including at moving grid points: the singular
pieces vanish to the required orders and the regular pieces match across
each periodic endpoint, so the boundary terms cancel. Hence
$\partial_{\boldsymbol a}^{\boldsymbol r}\mathcal L$ is the
meromorphic integral of the product of the three argument-derivative
generators.

Near $x_0\in S_i$ the singular term of factor $i$ is
\[
 (-1)^{r_i}(1-\lambda_i)_{r_i}
 p_i^{\lambda_i-r_i-1}t^{\lambda_i-r_i-1}.
\]
The other two factors are smooth there.  Subtract their Taylor
polynomial through degree $r_i$.  The degree-$r_i$ term alone has
denominator $\lambda_i$ near the origin; summed over $S_i$ its
numerator is exactly \eqref{xid:eq:triple-local-B}.
As in Theorem~\ref{xid:thm:coordinate-monomial}, its meromorphic integral
over $(0,h)$ is $B_i h^{\lambda_i}/\lambda_i$, whereas the generating
function of its coordinate finite parts is
$B_i(h^{\lambda_i}-1)/\lambda_i$.  Their difference is the complete
quantity $B_i/\lambda_i$.

The lower Taylor degrees have denominators bounded away from zero
on a small spectral polydisc.  The remainders are locally normally
integrable there.  A partition of unity treats the finitely many
disjoint singular points independently.  Subtracting all three
complete numerators therefore leaves a holomorphic function whose
Taylor coefficients are exactly the stated finite-part integrals.
The factorials in \eqref{xid:eq:triple-FP-value} follow from the exponential
Stieltjes generating convention.
\end{proof}

For zero argument derivatives the local term simplifies to
\[
 B_i^{\boldsymbol0}(\boldsymbol\lambda)
 =p_i^{\lambda_i-1}\sum_{x_0\in S_i}
             \prod_{j\ne i}\mathscr G(\lambda_j,\eta_{ij}(x_0)).
\]
Even here its full dependence on $\lambda_i$ must be retained; using
only $B_i(0,\lambda_j,\lambda_k)$ omits the powers of $\log p_i$.
Thus the character calculation and the ambient-coordinate theorem are
separate necessary parts of the reduction.

\subsection{Independent polynomial checks}

At positive integer $\lambda_i$, the identity
$\zeta(1-\lambda_i,y)=-B_{\lambda_i}(y)/\lambda_i$ makes
\eqref{xid:eq:triple-K} a rational piecewise-polynomial integral whenever
the shifts are rational.  This provides a check independent of Fourier
or special-function quadrature.  For
\[
 (p_1,p_2,p_3)=(1,2,3),\qquad
 (a_1,a_2,a_3)=(0,1/5,1/7),
\]
the exact integrations give
\begin{align}
 K(2,2,2)&=-\frac{49960598339}{1715322420000000},\notag\\
 K(2,3,2)&=-\frac{1641301145317}{315190494675000000}.
 \label{xid:eq:triple-rational-checks}
\end{align}
The supplied program compares these exact rational values with the
independent filtered Mellin integrals \eqref{xid:eq:weighted-T-mellin}.
The spectral check does not replace the local-subtraction proof for
Stieltjes jets, which is given above.
